\chapter{Fall Risk Assessment}
\section*{What Is This Test?} Fall-risk assessment identifies factors that make a person likely to fall and helps prevent injury.
\section*{Alternative Names} Fall evaluation; falls assessment; gait and balance assessment.
\section*{Principle} History, medication review, vision, orthostatic measurements, gait, strength, balance, and functional tests estimate risk.
\section*{Why This Test Is Ordered} It is used after a fall, recurrent falls, dizziness, weakness, mobility change, or in people with high-risk illness or medicines.
\section*{Primary vs Secondary Test} Clinical assessment is primary; ECG, imaging, laboratory tests, or vestibular studies are selected for the suspected cause.
\section*{How This Test Is Performed} The clinician reviews prior falls and observes standing, walking, turning, transfers, and use of aids.
\section*{What Preparation Is Needed} Bring medicine and fall history and use usual glasses, shoes, and walking aids.
\section*{Understanding Results} Risk increases with impaired balance, low blood pressure on standing, vision loss, unsafe environment, sedating medicines, or weakness.
\section*{Follow-up Tests} Medication review, physical therapy, vision assessment, ECG, orthostatic testing, or home-safety evaluation may follow.
\section*{References} \begin{enumerate}\item MedlinePlus. Fall Risk Assessment.\item Centers for Disease Control and Prevention. STEADI fall-prevention resources.\end{enumerate}
\clearpage\section*{Understanding Results and Follow-up}
The laboratory report should identify the specimen, method, units, reference interval or decision limit, and whether the result is positive, negative, abnormal, or inconclusive. Reference intervals vary with age, sex, pregnancy, specimen type, assay, and laboratory; a result outside a range is not automatically a diagnosis, and a result within a range does not always exclude disease. Discuss unexpected results with the ordering clinician, who can decide whether repeat or confirmatory testing, treatment, or referral is appropriate. Do not change medicines or delay urgent care based on an isolated result.
\section*{Regulatory Status and Authorization Date}This chapter describes a test category, not a specific commercial product. FDA, CDSCO, EU IVDR, and MHRA approval or registration dates therefore cannot be assigned without the assay manufacturer, product name, instrument, and intended use. For laboratory-developed tests, an individual agency approval date may not exist. See the Preface for the interpretation of regulatory status.
\clearpage
